% GENERATED FILE. Edit canonical lessons, then run npm run generate:books.
% canonical-source-hash: fnv1a64:2e45c65a5186b472
% canonical-lessons: KA-S08-vowel-sign-aa, KA-C08-dayavittu, KA-S117-letter-ya

\chapter{Please}
\label{ch:ka-please}
\hlchaptermodality{8}

\begin{chapteropening}
\textbf{By the end of this chapter:} \emph{I can say please in Kannada, and use the polite verb ending Kannada reaches for far more often than the word itself.}

Build \kn{ದಯವಿಟ್ಟು} out of daya (\textquotedblleft{}compassion\textquotedblright{}) + iṭṭu (\textquotedblleft{}having placed\textquotedblright{}), then make the everyday polite request instead: \kn{ಕುಳಿತುಕೊಳ್ಳಿ}, \textquotedblleft{}please sit.\textquotedblright{}
\end{chapteropening}

\section[ā]{\kn{◌ಾ} — one character, met inside words you already say}

\label{lesson:KA-S08-vowel-sign-aa}



\begin{quote}
Before the new one: \kn{ರ} — what does it do?

One character this time. Just one — and you have been saying it for pages without knowing which mark on the page it was.
\end{quote}

\subsection*{Script you'll notice: \kn{◌ಾ}}
\textbf{\kn{◌ಾ}} — \emph{ā}.

It is a \textbf{vowel sign}. It is not a letter and never stands alone: it attaches to a consonant and \textbf{replaces} the \emph{a} built into it with \emph{ā}. Replaces, not adds — the \emph{a} is gone.

You already say these, and every one of them has \kn{◌ಾ} somewhere inside it:

\begin{itemize}
\raggedright
  \item \textbf{\kn{ನಮಸ್ಕಾರ}} \emph{namaskāra} — hello / greetings (namaskāra — \textquotedblleft{}a making of a bow\textquotedblright{})
  \item \textbf{\kn{ಧನ್ಯವಾದ}} \emph{dhanyavāda} — thank you (dhanyavāda — \textquotedblleft{}an utterance of 'worthy'\textquotedblright{})
  \item \textbf{\kn{ನೀವು} \kn{ಹೇಗಿದ್ದೀರಾ}?} \emph{nīvu hēgiddīrā} — how are you? (respectful)
  \item \textbf{\kn{ನಾನು}} \emph{nānu} — I
\end{itemize}

\subsection*{Writing: \kn{◌ಾ} — copy what you see}
Put your pen on \kn{◌ಾ} and follow its line. Copy the shape you can see — slowly, and larger than it is printed.

\begin{quote}
This book does not yet tell you \textbf{where to start the character or which way to travel}. That is a real thing, taught with real variation from school to school, and it is not written down here until it can be written down with a source. Copying what is in front of you needs no such source, and it is how the shape gets into your hand in the meantime.
\end{quote}

\subsection*{Your turn}
\begin{itemize}
\raggedright
  \item \emph{Look:} at these words, and find \kn{◌ಾ} in the ones that have it
\end{itemize}

\begin{quote}
\kn{ನಮಸ್ಕಾರ}  ·  \kn{ಧನ್ಯವಾದ}  ·  \kn{ಹೌದು}
\end{quote}

\begin{itemize}
\raggedright
  \item \emph{Trace it:} \kn{◌ಾ} three times, saying \emph{ā} as you finish each one
  \item \emph{Look:} back at any page of this chapter and find \kn{◌ಾ} once more
\end{itemize}

\begin{tcolorbox}[breakable,colback=teal!4,colframe=teal!35!black,title={Before you move on}]
Which character is this — \kn{◌ಾ}? What vowel does it put on the consonant it attaches to — and what does it take off? (\textbf{Puts \emph{ā} on; takes the built-in \emph{a} off.}) Name one word you already say that contains it.
\end{tcolorbox}
\section[dayaviṭṭu]{\kn{ದಯವಿಟ್ಟು} (dayaviṭṭu) — please (having placed compassion)}

\label{lesson:KA-C08-dayavittu}



\begin{quote}
Kannada's everyday \textquotedblleft{}please\textquotedblright{} is, once you open it up, a small act of generosity: \textquotedblleft{}having \textbf{placed} compassion\textquotedblright{} between us.
\end{quote}

\begin{cousinweb}
\textbf{\kn{ದಯವಿಟ್ಟು}} = \textbf{please}, literally \textquotedblleft{}\textbf{having placed compassion / kindness}.\textquotedblright{} Two pieces:

\begin{itemize}
\raggedright
  \item \textbf{\kn{ದಯ}} (\emph{daya}) = \textbf{compassion, kindness, grace} — from Sanskrit \textbf{daya}.
  \item \textbf{\kn{ಇಟ್ಟು}} (\emph{iṭṭu}) = \textquotedblleft{}\textbf{having placed / set down},\textquotedblright{} from the verb \textbf{\kn{ಇಡು}} (\emph{iḍu}) \textquotedblleft{}to place, to put.\textquotedblright{}
\end{itemize}

So \emph{dayaviṭṭu} means, at heart, \emph{setting kindness down (before you)} — asking the listener to act from compassion.

That is exactly how \textbf{Tamil} builds it (\emph{tayavu seytu}, \textquotedblleft{}\textbf{do} the kindness\textquotedblright{}), and the same instinct as \textbf{Arabic} \emph{min faḍlik} (\textquotedblleft{}from your \textbf{grace}\textquotedblright{}) and \textbf{Hindi} \emph{kṛpayā} (from \emph{kṛpā}, \textquotedblleft{}\textbf{compassion}\textquotedblright{}). A wide arc of Asia says \emph{please} by naming the kindness it hopes for.
\end{cousinweb}

\begin{cousinweb}
\textbf{daya} + a light verb, four times over:

\begin{itemize}
\raggedright
  \item Kannada \textbf{\kn{ದಯವಿಟ್ಟು}} \emph{daya\textbf{viṭṭu}} — compassion \textbf{+ having placed}
  \item Tamil \textbf{\ta{தயவுசெய்து}} \emph{tayavu \textbf{sey}tu} — kindness \textbf{+ do}
  \item Telugu \textbf{\te{దయచేసి}} \emph{daya\textbf{cēsi}} — compassion \textbf{+ having done}
  \item Malayalam \textbf{\ml{ദയവായി}} \emph{daya\textbf{vāyi}} — compassion \textbf{+ as / having become}
\end{itemize}
\end{cousinweb}

\begin{culture}
The honest note: a standalone please-word is a little \textbf{formal / emphatic}. Everyday Kannada politeness leans on the \textbf{respectful command} — the verb ending \textbf{‑\kn{ಿ}} (\emph{-i}): \textbf{\kn{ಕುಳಿತುಕೊಳ್ಳಿ}} (\emph{kuḷitukoḷḷi}) \textquotedblleft{}please sit,\textquotedblright{}  \textbf{\kn{ಹೇಳಿ}} (\emph{hēḷi}) \textquotedblleft{}please tell me.\textquotedblright{} That ending already means \textquotedblleft{}please\textquotedblright{} do this. Add \textbf{\kn{ದಯವಿಟ್ಟು}} in front when you want to underline the courtesy.
\end{culture}

\begin{sounds}
Look at \textbf{\kn{ಟ್ಟು}} — a \textbf{doubled} \emph{ṭ}. The first \textbf{\kn{ಟ್}} is \emph{ṭa} silenced by the \textbf{virama} (its vowel cut), and it tucks \textbf{under} the next letter as an \emph{ottu} (a subscript form) before \textbf{\kn{ಟು}} \emph{ṭu}. Kannada writes a double consonant by stacking, not repeating side by side — a neat contrast with how the Roman alphabet just writes \textquotedblleft{}tt.\textquotedblright{}
\end{sounds}

\subsection*{Your turn}
Say these aloud:

\begin{itemize}
\raggedright
  \item \textquotedblleft{}daya\textquotedblright{} — compassion — then \textquotedblleft{}iṭṭu,\textquotedblright{} having placed
  \item the whole word — \textquotedblleft{}dayaviṭṭu\textquotedblright{} = please
  \item the everyday way — the verb ending: \textquotedblleft{}kuḷitukoḷḷi\textquotedblright{} = please sit
\end{itemize}

\begin{tcolorbox}[breakable,colback=teal!4,colframe=teal!35!black,title={Before you move on}]
What is the Kannada word for please? (\textbf{\kn{ದಯವಿಟ್ಟು}} \emph{dayaviṭṭu}.) What do its pieces mean? (\textquotedblleft{}\textbf{Compassion}\textquotedblright{} \emph{daya} + \textquotedblleft{}\textbf{having placed}\textquotedblright{} \emph{iṭṭu}.) Which cousins ask the same way? (\textbf{Tamil, Telugu, Malayalam} with \emph{daya} — plus \textbf{Arabic} \emph{faḍl}, \textbf{Hindi} \emph{kṛpā}.) What carries \textquotedblleft{}please\textquotedblright{} in ordinary speech? (\textbf{The respectful verb ending ‑\kn{ಿ}} \emph{-i}.)
\end{tcolorbox}
\section[ya]{\kn{ಯ} — one character, met inside words you already say}

\label{lesson:KA-S117-letter-ya}



\begin{quote}
Before the new one: \kn{ಕ} — what does it do?

One character this time. Just one — and you have been saying it for pages without knowing which mark on the page it was.
\end{quote}

\subsection*{Script you'll notice: \kn{ಯ}}
\textbf{\kn{ಯ}} — \emph{ya}.

It is a \textbf{consonant}, and in this script a consonant is never bare: it comes with an \emph{a} already in it. So it is not \emph{y}, it is \textbf{ya}.

You already say these, and every one of them has \kn{ಯ} somewhere inside it:

\begin{itemize}
\raggedright
  \item \textbf{\kn{ಧನ್ಯವಾದ}} \emph{dhanyavāda} — thank you (dhanyavāda — \textquotedblleft{}an utterance of 'worthy'\textquotedblright{})
  \item \textbf{\kn{ದಯವಿಟ್ಟು}} \emph{dayaviṭṭu} — please (dayaviṭṭu
\end{itemize}

\subsection*{Writing: \kn{ಯ} — copy what you see}
Put your pen on \kn{ಯ} and follow its line. Copy the shape you can see — slowly, and larger than it is printed.

\begin{quote}
This book does not yet tell you \textbf{where to start the character or which way to travel}. That is a real thing, taught with real variation from school to school, and it is not written down here until it can be written down with a source. Copying what is in front of you needs no such source, and it is how the shape gets into your hand in the meantime.
\end{quote}

\subsection*{Your turn}
\begin{itemize}
\raggedright
  \item \emph{Look:} at these words, and find \kn{ಯ} in the ones that have it
\end{itemize}

\begin{quote}
\kn{ಧನ್ಯವಾದ}  ·  \kn{ದಯವಿಟ್ಟು}  ·  \kn{ನಮಸ್ಕಾರ}
\end{quote}

\begin{itemize}
\raggedright
  \item \emph{Trace it:} \kn{ಯ} three times, saying \emph{ya} as you finish each one
  \item \emph{Look:} back at any page of this chapter and find \kn{ಯ} once more
\end{itemize}

\begin{tcolorbox}[breakable,colback=teal!4,colframe=teal!35!black,title={Before you move on}]
Which character is this — \kn{ಯ}? What sound does it carry? (\emph{\textbf{ya}}.) Name one word you already say that contains it.
\end{tcolorbox}
